\documentclass[a4paper]{article}

\usepackage[a4paper, top=8mm, bottom=8mm, left=8mm, right=8mm]{geometry}

\usepackage{polyglossia}
\setdefaultlanguage[babelshorthands=true]{russian}
\setotherlanguage{english}

\usepackage{fontspec}
\setmainfont{FreeSerif}
\newfontfamily{\russianfonttt}[Scale=0.7]{DejaVuSansMono}

\usepackage[tiny, compact]{titlesec}

\usepackage{titling}
\setlength{\droptitle}{-1cm}
\pretitle{\begin{center}\begin{bfseries}\Large}
\posttitle{\par\end{bfseries}\end{center}}
\preauthor{\begin{center}\normalsize}
\postauthor{\par\end{center}\vspace{-1.8cm}}

\usepackage{hyperref}
\usepackage{bookmark}
\usepackage{csquotes}
\usepackage{graphicx}
\usepackage{amsmath}
\usepackage{booktabs}
\usepackage{siunitx}
\usepackage{float}
\restylefloat{figure}

\title{Исследование влияния иерархического кэширования на производительность матричных вычислений в OpenCL на архитектурах с унифицированной памятью}

\author{Симонов Кирилл Андреевич}

\date{}

\begin{document}

\maketitle

\begin{flushright}
\begin{tabular}{r}
    Группа: \emph{\enquote{25.Б41-мм}} \\
    Кафедра: \emph{математического обеспечения и администрирования информационных систем} \\
    Научный руководитель: \emph{Григорьев Семён Вячеславович} \\
    Номер семестра практики: \emph{2} \\
\end{tabular}
\end{flushright}

\section{Постановка задачи}

Выполнение операции перемножения плотных матриц (SGEMM) лежит в основе подавляющего большинства алгоритмов машинного обучения и вычислительной математики. Ввиду кубической асимптотической сложности данной задачи, ключевым фактором, ограничивающим скорость ее выполнения на графических ускорителях (GPU), становится не максимальная производительность арифметико-логических устройств, а пропускная способность шины памяти.

Данная учебная практика посвящена исследованию подходов к оптимизации доступа к памяти при вычислении матричных произведений. В отличие от дискретных видеокарт, целевой платформой в работе выступает система на кристалле (SoC) с архитектурой унифицированной памяти (Unified Memory Architecture, UMA). В таких системах центральный и графический процессоры физически разделяют общий объем оперативной памяти, что требует специфического подхода к распределению нагрузок и утилизации локальной кеш-памяти. Основой для проведения экспериментов послужил проект myGEMM, предоставляющий набор последовательно усложняющихся реализаций алгоритма.

Особенностью macOS является использование транслятора \texttt{cl2Metal}, который преобразует OpenCL C в промежуточное представление Metal. Оригинальный код myGEMM содержит сложные препроцессорные вычисления (например, \texttt{\#define RTS (TS/WPT)}), что приводит к аварийному завершению компиляции под \texttt{cl2Metal}. Таким образом, требуется не только оптимизация самих алгоритмов, но и адаптация механизма передачи параметров.

\textbf{Целью работы является} адаптация существующих алгоритмов блочного матричного умножения под трансляторы macOS, а также экспериментальная оценка влияния двумерного блочного разбиения и регистровой оптимизации на итоговую производительность в условиях UMA.

Для выполнения поставленной цели потребовалось решить следующие задачи:
\begin{itemize}
    \item модифицировать кодовую базу ядер проекта myGEMM\footnote{Cedric Nugteren. Репозиторий проекта myGEMM. URL: \url{https://github.com/CNugteren/myGEMM} (дата обращения: 21.03.2026).} для обеспечения совместимости с компилятором \texttt{cl2Metal};
    \item спроектировать автономную систему программ для статической генерации параметров разбиения и управления запусками;
    \item провести анализ производительности различных конфигураций ядер на целевом кристалле и выявить наиболее эффективные стратегии работы с памятью.
\end{itemize}

\section{Описание предлагаемого решения}

Оригинальная архитектура проекта myGEMM опирается на передачу параметров оптимизации (размеров блоков, ширины векторизации) непосредственно на этапе JIT-компиляции OpenCL-кода через динамически задаваемые параметры препроцессора. Однако встроенный в подсистему macOS компилятор \texttt{cl2Metal} демонстрирует нестабильную работу при разборе сложных препроцессорных выражений внутри \texttt{.cl} файлов, что делает невозможным использование оригинального подхода на платформе Apple Silicon.

Для решения этой проблемы был разработан независимый программный комплекс на языке Python, выполняющий расчет параметров оптимизации на центральном процессоре. Инструментарий состоит из следующих изолированных компонентов:

\begin{itemize}
    \item \texttt{calculate\_kernels.py} — модуль статической подготовки конфигураций. Указанная программа служит инструментом для развертывания пользовательских параметров оптимизации (\texttt{TS}, \texttt{WPT}, \texttt{WIDTH} и других) в виде фиксированных констант препроцессора. Программа выполняет вычисление производных величин, необходимых для работы алгоритмов (например, \texttt{RTS = TS/WPT}, \texttt{LPT = (TSDK * WPT) / TS}), и записывает их в заголовочный файл \texttt{settings.h} как готовые целочисленные значения. Это полностью переносит математическую нагрузку с препроцессора OpenCL на сторону Python, что необходимо для стабильной трансляции кода в \texttt{cl2Metal} на платформе Apple Silicon. Кроме того, модуль обеспечивает реализацию векторизации: в зависимости от заданного параметра \texttt{WIDTH}, автоматически генерируется директива \texttt{typedef}, определяющая тип \texttt{floatX} (\texttt{float}, \texttt{float2} или \texttt{float4}), а также набор вспомогательных параметров препроцессора, что позволяет коду ядер OpenCL единообразно работать с векторами данных, сохраняя совместимость с управляющей программой на C++.
    \item \texttt{analyze\_results.py} — аналитический модуль, предназначенный для обработки данных по завершении сессии экспериментов. Программный модуль осуществляет синтаксический анализ исходных файлов результатов, рассчитывает производительность в GFLOPS для каждой протестированной конфигурации и выполняет статистическую обработку: расчет среднего арифметического, вычисление 95\% доверительных интервалов и выявление статистических выбросов. Это позволяет отсеивать результаты, искаженные фоновой нагрузкой системы, и выявлять наиболее производительную конфигурацию вычислительного ядра для конкретной аппаратной платформы.
\end{itemize}

В целях обеспечения отказоустойчивости разработанного программного комплекса внедрены модульные тесты на базе среды тестирования \texttt{pytest}. Тесты проверяют алгоритмы отбраковки выбросов и корректность расчета статистических показателей. Контроль качества кода реализован на базе платформы непрерывной интеграции GitHub Actions.

\section{Экспериментальное исследование производительности}

\subsection{Условия проведения эксперимента}

Оценка эффективности реализованных алгоритмов проводилась на изолированном стенде, конфигурация которого минимизирует влияние фоновых процессов ОС:
\begin{itemize}
    \item кристалл Apple M3 (8 CPU, 10 GPU-ядер, аппаратное динамическое кэширование);
    \item 8 ГБ унифицированной памяти (Unified Memory);
    \item операционная система macOS Tahoe 26.1;
    \item графический стек OpenCL 1.2 поверх Metal API;
    \item компилятор Apple Clang с флагами оптимизации \texttt{-O3 -Wall -framework OpenCL}.
\end{itemize}

В качестве входных данных использовались квадратные матрицы размерностью \(4096 \times 4096\) и \(8192 \times 8192\) элементов. Для каждого ядра выполнялось 50 контрольных итераций, которым предшествовали 15 предварительных запусков для стабилизации температурного режима кристалла и предотвращения динамического снижения тактовой частоты. Между сменой тестируемых размерностей выдерживалась трехсекундная программная пауза.

Для минимизации влияния фоновой активности операционной системы macOS (обусловленной работой таких процессов, как \texttt{WindowServer} или \texttt{Spotlight}) на точность замеров, в методику проведения эксперимента был включен этап предварительной стабилизации графического процессора (\textit{warm-up}). Перед началом записи данных выполнялось 15 инициализирующих итераций, что обеспечивало выравнивание тактовой частоты графического процессора и выход системы из энергосберегающих состояний. Все 50 итераций измерения включались в финальный расчет, что при достаточном объеме выборки позволяет компенсировать случайные кратковременные всплески времени выполнения, связанные с планировщиком задач ОС. Стабильность полученных результатов при различных нагрузках подтверждается высокой повторяемостью средних значений производительности (GFLOPS) при последующих сессиях экспериментов.

Оценка вычислительной производительности основывалась на подсчете затраченного физического времени \(t\). Поскольку для вычисления одного элемента результирующей матрицы \(N \times N\) требуется \(2N\) скалярных операций сложения и умножения, итоговый показатель производительности в GFLOPS рассчитывался по формуле:
\[
GFLOPS = \frac{2N^3}{t \cdot 10^9}.
\]

\subsection{Результаты анализа производительности ядер}

Сводные результаты тестирования производительности для различных размерностей матриц представлены на рисунке~\ref{fig:kernel-diagrams}.

\vspace{-2cm}
\begin{figure}[H]
    \centering
    \begin{minipage}[t]{0.45\textwidth}
        \centering
        \includegraphics[width=\textwidth]{diagrams/kernel_4096.pdf}
        \\[-2cm] \centering \textbf{а)} Размер матрицы $4096 \times 4096$
    \end{minipage}
    \hfill
    \begin{minipage}[t]{0.45\textwidth}
        \centering
        \includegraphics[width=\textwidth]{diagrams/kernel_8192.pdf}
        \\[-2cm] \centering \textbf{б)} Размер матрицы $8192 \times 8192$
    \end{minipage}
    \vspace{-1.5cm}
    \caption{Динамика производительности вычислительных ядер в зависимости от стратегии работы с памятью}
    \label{fig:kernel-diagrams}
\end{figure}

Численные значения средней производительности (GFLOPS) для матрицы \(4096\times4096\): Kernel~1 — 160, Kernel~4 — 313, Kernel~8 — 836. Для матрицы \(8192\times8192\): Kernel~1 — 147, Kernel~4 — 283, Kernel~8 — 791. Таким образом, комбинация двумерного блочного разбиения и регистровой оптимизации (\text{Kernel~8}) обеспечивает увеличение производительности в 5,2–5,7 раз в сравнении с базовой реализацией для двух масштабов матриц.

Детальный анализ полученных показателей выявил выраженную корреляцию между степенью локализации данных и итоговой производительностью на UMA-архитектуре:

\begin{itemize}
    \item \textbf{Базовое ядро (Kernel 1)} продемонстрировало наименьшую эффективность (около 120–160 GFLOPS) ввиду высокой конкуренции потоков за доступ к глобальной шине памяти;
    \item \textbf{Векторизация чтения (Kernel 4)} увеличила пропускную способность за счет использования типов \texttt{float4}, дав прирост производительности примерно в 3,5 раза в сравнении с базовой версией;
    \item \textbf{Комбинация двумерного блочного разбиения и регистровой оптимизации} (\text{Kernel 8}) позволила достичь максимальной производительности (более 750 GFLOPS). Загрузка блоков матриц в локальную память общего доступа и аккумуляция промежуточных сумм в регистрах практически полностью изолировали вычислительные конвейеры GPU от основной системной памяти.
\end{itemize}

Для проверки исследовательских вопросов были сопоставлены ядра с близкими оптимизациями. Установлено, что добавление регистровой оптимизации к двумерному блочному разбиению (\text{Kernel~7} против \text{Kernel~6}) обеспечивает дополнительное увеличение производительности на 9\%, а замена векторизации с \texttt{float4} на \texttt{float2} (\text{Kernel~8} против \text{Kernel~7}) снижает производительность на 5\%, что свидетельствует о предпочтительности более широких векторов на Apple M3.

Дополнительная проверка распределения временных отклонений с помощью критерия Шапиро — Уилка подтвердила, что для высокооптимизированного ядра (\text{Kernel 8}) случайные задержки планировщика macOS перестают оказывать значимое влияние на работу GPU, переводя задачу из состояния ограничения пропускной способностью памяти в стабильный режим ограничения вычислительной мощностью.

\section{Заключение}

В ходе выполнения практики были успешно реализованы поставленные задачи:
\begin{itemize}
    \item устранены архитектурные ограничения компилятора \texttt{cl2Metal} путем разработки Python-комплекса (\texttt{calculate\_kernels.py}) для статической генерации заголовочных файлов; все параметры блочного разбиения вычисляются на уровне центрального процессора, что исключает необходимость выполнения вычислений на этапе препроцессинга в ядрах OpenCL.
    \item проведены эксперименты на чипе Apple M3, подтвердившие критическую важность двумерного блочного разбиения для платформ с UMA-архитектурой; получены количественные оценки: Kernel~8 (двумерное блочное разбиние и регистровая оптимизация) превосходит Kernel~1 в 5,2–5,7 раз, а Kernel~4 (только векторизация) — в 2,3–2,5 раза.
    \item определена конфигурация ядра (Kernel 8), позволившая достичь пиковых 836 GFLOPS на матрице \(4096\times4096\), что более чем в 5 раз превосходит показатели наивной реализации.
\end{itemize}

Исходный код разработанного инструментария, конфигурационные файлы для CI/CD размещены в открытом доступе. Структура репозитория и описание конфигураций позволяют воспроизвести эксперименты на аналогичном аппаратном обеспечении.

Репозиторий проекта: \url{https://github.com/NewsPGG/myGEMM}.

\end{document}